\section{Evaluation}
\label{sec:evaluation}

This section reports the system's measured performance: the dataset the models were trained and
evaluated on, the classification accuracy of the on-device cascade on a held-out split, the behaviour of
the gated audio confirmation stage, and the on-device resource footprint. Accuracy and confusion are
measured on a held-out test split of the collar's own recorded clips; the resource footprint is read from
the firmware build. On-hardware timing (inference latency and over-the-air transfer) and power draw are
not yet instrumented and are noted as future work (\cref{sec:future-work}).

\subsection{Dataset}

The models were trained and evaluated on labelled behaviour clips recorded from a single cat with the
collar (\cref{sec:training-pipeline}), each a 5\,s window of $\mu$-law audio paired with time-aligned
motion. Of 191 recorded clips, 155 carry a behaviour label; the remaining 36 are unlabelled and are
excluded from training and evaluation. The label set is discovered from the data rather than fixed
(\cref{sec:on-device-ai}); on this recording it is \emph{eat}, \emph{drink}, \emph{resting} and
\emph{moving}.

\begin{table}[ht]
  \centering
  \caption{Dataset summary.}
  \label{tab:dataset}
  \begin{tabularx}{\textwidth}{@{}lX@{}}
    \toprule
    \textbf{Property} & \textbf{Value} \\
    \midrule
    Cats                      & 1 \\
    Capture sessions          & 19 \\
    Clip length               & 5\,s \\
    Audio sample rate         & 8\,kHz ($\mu$-law) \\
    IMU sample rate           & 104\,Hz (6-axis) \\
    Labelled clips            & 155 (a further 36 are unlabelled and excluded) \\
    Train / test split        & 116 / 39 clips (held-out 25\,\%, stratified by class) \\
    \bottomrule
  \end{tabularx}
\end{table}

\begin{table}[ht]
  \centering
  \caption{Labelled clips per behaviour class.}
  \label{tab:dataset-classes}
  \begin{tabularx}{\textwidth}{@{}Xr@{}}
    \toprule
    \textbf{Class} & \textbf{Clips} \\
    \midrule
    Eat      & 94 \\
    Resting  & 40 \\
    Drink    & 11 \\
    Moving   & 10 \\
    \midrule
    \textbf{Total (labelled)} & 155 \\
    \bottomrule
  \end{tabularx}
\end{table}

\subsection{Classification accuracy}

The collar runs the confidence-gated cascade of \cref{sec:on-device-ai}: the IMU stage classifies every
window, and the audio stage confirms when the IMU result is uncertain (softmax confidence below the
operating threshold of 0.75). On the held-out test split the cascade reaches an overall accuracy of
\textbf{92.3\,\%} (36 of 39 clips). On this dataset the IMU stage alone reaches the same 92.3\,\%: the
gated audio stage fired but did not change the outcome (\cref{tab:confusion} and the following
subsection).

\begin{table}[ht]
  \centering
  \caption{Per-class precision / recall / F1 on the held-out split. On this data the full cascade is
           identical to the IMU stage alone, so a single set of figures is shown.}
  \label{tab:accuracy}
  \begin{tabularx}{\textwidth}{@{}X rrr@{}}
    \toprule
    \textbf{Class} & \textbf{Precision} & \textbf{Recall} & \textbf{F1} \\
    \midrule
    Eat      & 0.89 & 1.00 & 0.94 \\
    Resting  & 1.00 & 1.00 & 1.00 \\
    Drink    & 1.00 & 0.67 & 0.80 \\
    Moving   & 0.00 & 0.00 & 0.00 \\
    \midrule
    \textbf{Overall accuracy} & \multicolumn{3}{c}{92.3\,\%} \\
    \bottomrule
  \end{tabularx}
\end{table}

\begin{table}[ht]
  \centering
  \caption{Confusion matrix on the held-out split (rows = actual, columns = predicted).}
  \label{tab:confusion}
  \begin{tabular}{@{}l rrrr@{}}
    \toprule
    & \multicolumn{4}{c}{\textbf{Predicted}} \\
    \cmidrule(lr){2-5}
    \textbf{Actual} & Drink & Eat & Moving & Resting \\
    \midrule
    Drink   & 2 & 1  & 0 & 0 \\
    Eat     & 0 & 24 & 0 & 0 \\
    Moving  & 0 & 2  & 0 & 0 \\
    Resting & 0 & 0  & 0 & 10 \\
    \bottomrule
  \end{tabular}
\end{table}

\subsection{The gated audio confirmation stage}

The audio stage exists to separate behaviours that look alike to an accelerometer but differ by sound,
eating and drinking being the canonical case. On the held-out split it fired on 6 of the 39 test clips
(15\,\%), the windows where the IMU stage sat below the 0.75 confidence threshold. On this small
single-cat set it changed none of those predictions and corrected no errors, so the cascade's accuracy
equals the IMU stage's. The interpretation is not that the stage is redundant but that it was not needed
here: the IMU model was already confident and correct on the confusable cases present in this recording,
so the gated audio stage had nothing to overturn. Establishing a measured accuracy gain from the audio
stage calls for a larger and more behaviourally ambiguous dataset than this first single-cat recording,
and is left to future work (\cref{sec:future-work}). The value demonstrated here is the mechanism itself:
a second, independent sensing modality that is consulted only when the primary stage is unsure, at no
extra cost when it is confident.

\subsection{Resource footprint}

The collar is the binding resource constraint. The tensor arena is statically allocated and dominates the
RAM cost of inference; the audio window buffer and the BLE stack (now including LE Secure Connections
pairing) make up most of the rest.

\begin{table}[ht]
  \centering
  \caption{Collar resource footprint (the constraining device), read from the firmware build.}
  \label{tab:footprint}
  \begin{tabularx}{\textwidth}{@{}llX@{}}
    \toprule
    \textbf{Resource} & \textbf{Value} & \textbf{Notes} \\
    \midrule
    RAM (static)        & 226\,KiB of 256\,KiB & Static \texttt{data+bss} with inference and the encrypted
                                                 BLE link enabled (88.3\,\%); includes the shared tensor
                                                 arena. 231\,364\,B \\
    Flash (image)       & 337\,KiB             & Application image (\texttt{text+data}), excluding the
                                                 model partition (63.4\,\%). 345\,476\,B \\
    Shared tensor arena & 48\,KiB              & One arena for both stages; only one model is resident at a
                                                 time (IMU and audio flip in and out) \\
    Model storage       & 256\,KiB flash       & Two OTA slots; the interpreter reads weights directly from
                                                 this partition (XIP), with only the shared arena in RAM \\
    IMU model size      & 23\,736\,B ($\approx$23\,KiB) & int8 \texttt{.tflite}; delivered to slot~0 over
                                                 the air, not in the firmware image \\
    Audio model size    & 19\,760\,B ($\approx$19\,KiB) & int8 \texttt{.tflite}; delivered to slot~1 \\
    \bottomrule
  \end{tabularx}
\end{table}

\subsection{Testing and verification}

The system's logic is covered by automated unit tests that run in continuous integration on every change,
and the integration and on-device behaviour is verified on the hardware. The security controls and their
evidence are detailed separately (\cref{sec:security-privacy}).

\begin{itemize}[leftmargin=1.5em]
  \item \textbf{Automated unit tests.} 80 Python tests cover the dashboard data layer and the Cloud
        Function's signal logic (per-window normalisation, windowing, WAV parsing, and the path-validation
        regexes), including property and fuzz tests that hammer the untrusted-input surface. Thirteen
        host-compiled C tests cover the on-device cascade (the rules-based activity gate and the
        confidence-gated IMU/audio gating) and the $\mu$-law audio codec and model-input representation,
        the latter swept exhaustively over the whole int16 domain and checked to match the training
        representation. The full per-test inventory is listed in \cref{app:tests}.
  \item \textbf{Continuous integration.} Every push and pull request runs the Python and C test suites and
        CodeQL static security analysis, with the workflow actions pinned to commit hashes.
  \item \textbf{On-hardware verification.} Model load-from-flash, on-device classification, the activity
        gate's rest/wake cycle, the encrypted collar-to-station link, and the build and backend-deploy
        steps were each confirmed directly on the collar and station.
\end{itemize}

\subsection{Limitations}

The evaluation is a first, single-cat, within-subject result and should be read as such. The dataset is
small (155 labelled clips) and imbalanced: eat dominates (94 clips) while moving and drink have only 10
and 11 clips, so their per-class figures are unreliable, moving in particular scoring zero on a two-clip
test slice. The held-out split is a single stratified 25\,\% partition, not cross-validated. The gated
audio stage produced no measured accuracy gain on this set (above); whether it helps on confusable
behaviours across more cats is the central question for future work (\cref{sec:future-work}). Finally,
on-device inference latency, over-the-air transfer time and power draw are not yet instrumented on the
hardware; the resource footprint above is the measured on-device figure available now.
